% Tapis de Quillen, with Quillen's offprint and preprint (cote 162-5, typed part): archive cote 162-5, site page 12 (PDF page 13).
% typescript. Draft compared: none.
% AI-transcribed by Claude Opus 5.5 (claude-opus-5-5, adaptive thinking, effort high), Message Batches API, batch t162_full (msgbatch_01UAh3zyWPd454AeDkhnmQ33), 2026-11-13; not yet checked by eye.
\scanpage{162-5/12}
\mspage{11}

III Point de vue ``motivique'' en théorie du cobordisme.

% NOTE: the right edge of the scan cuts off the ends of several typed lines; the visible text is transcribed as it stands
Soit C la catégorie des variétés différebtiables\inserted{(\unsure{compactes?})} (pas nécessairem
orinhtables), \inserted[typed]{\uline{les morphismes étant les classes d'homotopie d'applications continues}} Si B est une catégorie, on s'intéresse aux couples

$(F.,F^{\cdot})$ d'un foncteur covariant et d'un foncteur contravariant de C

dans B, satisfaisant les conditions que pour tout $X\in \mathrm{Ob}\, C$, on a $F.(X)=$

$F^{\cdot}(X)$, et que si on a un produit fibré ordinaire
\diagram{\begin{tikzcd}
 & X' \arrow[dl,"g'"'] \arrow[dr,"f'"] & \\
X \arrow[dr,"f"'] & & Y' \arrow[dl,"g"] \\
 & Y &
\end{tikzcd}}
,

avec f et g transversaux, alors on a $g^{\cdot}f. = f!\, g'^{\cdot}$\inserted[typed]{(on pose $f.=F.(f)$,$f^{\cdot}=F^{\cdot}($}. La question est

de trouver une B et un couple $(F.,F^{\cdot})$ qui soient ``universels'' dans u

sens évident. Quillen prouve que la construction suivante donne une

telle solution.

Pour toute variété $X\in \mathrm{ob}\, C$ , on définit un anneau commutatif

B(X) (l'anneau de bordisme de X) en prenant comme éléments les class

\xout{de} d'éléments de C au dessus de X, pour la relation de cobordisme:

$f'{:}X'\to X$ et $f''{:}X''\to X$ sont équivalents, s'il existe une variété à bord

Z, et un $f{:}Z\to X$, tel que bord $Z = X'\amalg X''$ et $f'$ et $f''$ sont induits

par f. La somme de B(X) est la somme disjointe (c'est une loi de grou

\xout{\unsure{inverse étant}} où tout élément est d'ordre \hand{\unsure{2}} \xout{\illegible{1}}2), le produit est

donné par le produit fibré, en utilisant le théorème de trahsversalit

de Thom pour bouger un peu f'', de sorte à le rendre transversal à f'.
\margin{NB L'anneau est de plus \unsure{gradué} par la dimension relative ....}
Si on a $f{:}\ X\to Y$, on définit $f. : B(X)\to B(Y)$ de façon évidente (tou

variété sur X est sur Y), et $f^{\cdot}{:}\ B(Y)\to B(X)$ par produit fibré (uti

sant encore le th. de transversalité). De cette façon, B(X) devient u

foncteur-anneau \inserted{(\unsure{gradué})} contravariant en X, et un foncteur groupe covariant e

\mspage{12}
